\section{Bounded components of distance slabs}
\label{sec:slabs}

The separating interfaces constructed in Section~\ref{sec:doubling-bubbles}
will now give a cover of each cut component by sets of bounded diameter.
The point is to find one interface component that every path from the
base point to a given slab component must cross.  Separate crossings of
unrelated interface components would not suffice for this purpose.

\begin{lemma}[The complementary-region tree]
\label{lem:region-tree}
Let $P$ be a connected smooth manifold, possibly with boundary.  Let
$F_1,\ldots,F_m$ be disjoint compact connected two-sided neatly embedded
hypersurfaces, each of which separates $P$.  The graph with one vertex for
each component of $P\setminus\bigcup_i F_i$ and one edge across each $F_i$
is a finite tree.  In particular, two connected subsets of the complement
that lie in different components are separated by some single $F_i$.
\end{lemma}

\begin{proof}
A bicollar of $F_i$ can be chosen disjoint from all the other
hypersurfaces.  Its two sides are connected, so each side belongs to a
unique complementary component.  This defines the endpoints of its
edge.  The complement has finitely many components: insert the
hypersurfaces one at a time, observing that removing a connected
two-sided hypersurface from a connected manifold creates at most two
components.  The same argument applies with boundary and to a neat
hypersurface, using its bicollar.

The graph is connected.  Indeed, points in complementary components
can first be moved into the interior of $P$ and then joined by an
interior path transverse to the hypersurfaces.  A transverse compact
path has finitely many crossings and hence gives a walk in the graph.
Every edge is a bridge: a graph path joining its endpoints without
that edge could be realized by paths in complementary regions and
bicollars, thereby joining the two sides of $F_i$ without meeting
$F_i$.  This contradicts separation by $F_i$.  A connected finite
graph in which every edge is a bridge is a tree.  Any edge on the
unique path between two distinct vertices gives the final assertion.
\end{proof}

Fix a connected component $P$ of the cut manifold of
Proposition~\ref{prop:separating-cut}.  As in
Section~\ref{sec:doubling-bubbles}, use the smoothed metric $h$ on its
double, let $d$ be the intrinsic length metric of $h|_P$, and fix
$o\in\operatorname{int}P$.  Write $r(x)=d(o,x)$.  The metric $d$ is
proper, and Lemma~\ref{lem:double-smoothing} gives
\begin{equation}
\label{eq:projection-distance}
 \dist_g(q(x),q(y))\le 2d(x,y) \qquad (x,y\in P).
\end{equation}
Here $\dist_g$ is distance in the original manifold $M$, not the
intrinsic distance in a cut component: project paths in $P$ and use
$g\le 4h$.

\begin{proposition}[Slab diameter bound]
\label{prop:slab-diameter}
For $\ell<u\le\ell+2$, every connected component $C$ of
\[
 \{x\in P:\ell<r(x)<u\}
\]
has $d$-diameter at most $220$.
\end{proposition}

\begin{proof}
If $\ell<100$, then $r<u<102$ on the slab.  The triangle inequality
through $o$ gives $\diam_d C\le 204$.

Suppose $\ell\ge100$ and the slab is nonempty.  Apply
Proposition~\ref{prop:separating-interfaces} at level $\ell$.  It gives
finitely many disjoint compact connected neat two-sided hypersurfaces
$F_1,\ldots,F_m$ in
\begin{equation}
\label{eq:interface-radius-window}
 \ell-91<r<\ell-18,
\end{equation}
each separating $P$ and each of $d$-diameter at most $4\pi$.
The full phase contains $o$, whereas the empty phase contains the
slab; the phase is constant on each complementary region.  Thus $o$
and $C$ lie in different complementary regions.  In this assertion,
$C$ lies in a single region because it is connected and misses every
$F_i$.

By Lemma~\ref{lem:region-tree}, some one component $F_i$ separates
$o$ from all of $C$.  Proper length spaces are geodesic, so for each
$x\in C$ choose a minimizing path from $o$ to $x$.  It meets $F_i$
at some point $z_x$.  Each subpath is minimizing, and hence its
remaining length after this crossing is
\begin{equation}
\label{eq:slab-residual-length}
 d(x,z_x)=r(x)-r(z_x)
 < (\ell+2)-(\ell-91)=93.
\end{equation}
For $x,y\in C$, the triangle inequality and the ambient
$d$-diameter bound for $F_i$ give
\[
 d(x,y)\le d(x,z_x)+d(z_x,z_y)+d(z_y,y)
 < 2\cdot93+4\pi<220.
\]
Together with the first case this proves the proposition.
\end{proof}

For the graph construction we use the particular open cover
\begin{equation}
\label{eq:slab-cover}
 \mathcal U_P=
 \left\{\text{connected components of }
   \{j-\tfrac34<r<j+\tfrac34\}:j=0,1,2,\ldots\right\}.
\end{equation}
These are open sets because a manifold with boundary is locally
path connected.  The intervals cover $[0,\infty)$, and no point
belongs to more than two of them.  Components for a fixed $j$ are
disjoint, so $\mathcal U_P$ has multiplicity at most two.  Each of
its sets has compact closure, since $r$ is proper.  By
Proposition~\ref{prop:slab-diameter} and
\eqref{eq:projection-distance},
\begin{equation}
\label{eq:cover-diameter}
 \diam_g q(U)\le440 \qquad(U\in\mathcal U_P).
\end{equation}
The cover itself need not be locally finite.  The next section makes
the partition of unity, and therefore the target graph, locally
finite without losing this diameter bound.
